% =========================================================================
% Chapter 5 - The VASP interface
% =========================================================================
\chapter{The VASP interface}
\label{ch:vasp}

Poraqu\^e computes the external ionic potential natively, and reproducing
VASP's own convention exactly required reading the Fortran source. This chapter
records what it says, because the details decide whether the Python
reconstruction agrees with a reference potential or misses it by $13\,\%$.

Reference \file{EXTCAR} and \file{TAUCAR} files were used during development to
\emph{validate} that reconstruction. They are not required to run Poraqu\^e ---
nothing in the pipeline reads them --- and the reference calculations
themselves are ordinary VASP~6.2.0 output in every other respect.

\section{The tags that produce the reference files}

Two \file{INCAR} flags, declared in \file{base.F} and parsed in \file{main.F}:

\begin{center}
\begin{tabular}{@{}lll@{}}
\toprule
Tag & Variable & Writes \\
\midrule
\opt{EXTCAR} & \opt{LEXTCAR} & local external (ionic) potential \\
\opt{TAUCAR} & \opt{LTAUCAR} & kinetic energy density \\
\bottomrule
\end{tabular}
\end{center}

\section{\texorpdfstring{\file{EXTCAR}}{EXTCAR}: the local pseudopotential}

\subsection{What is written}

\file{main.F} calls \opt{POTION} to build the potential in reciprocal space,
transforms it to real space, and writes it with \opt{OUTPOS} + \opt{OUTPOT}.
The choice of \opt{OUTPOT} rather than \opt{OUTCHG} is what makes \file{EXTCAR}
a \emph{potential} file in the \file{LOCPOT} sense: \textbf{no multiplication by
$\Omega$}. The source comment states the conventions directly: values in eV,
no volume scaling, $\GG=0$ set to zero, a single spin-independent block, no PAW
one-centre terms and no non-local projector contribution.

\subsection{The formula}

\opt{POTION} (\file{pot.F}) contains, per species and per $\GG$ vector:

\begin{lstlisting}[style=fortran]
ARGSC  = NPSPTS/P(NT)%PSGMAX
PSGMA2 = P(NT)%PSGMAX - P(NT)%PSGMAX/NPSPTS
ZZ     = -4*PI*P(NT)%ZVALF*FELECT

G = SQRT(GX**2+GY**2+GZ**2)*2*PI
IF ( (G /= 0) .AND. (G < PSGMA2) ) THEN
   I    = INT(G*ARGSC)+1
   REM  = G - P(NT)%PSP(I,1)
   VPST = PSP(I,2)+REM*(PSP(I,3)+REM*(PSP(I,4)+REM*PSP(I,5)))
   CVPS(N) = CVPS(N) + ( VPST + ZZ/G**2 ) / LATT_CUR%OMEGA * CSTRF(N,NT)
ELSE
   CVPS(N) = 0._q
ENDIF
\end{lstlisting}

which is
\begin{equation}
  \boxed{\;
  \vext(\GG) = \frac{1}{\Omega}\sum_{s} S_s(\GG)
    \left[\, v^{s}_{\mathrm{short}}(G)
      \;-\; \frac{4\pi Z^{\mathrm{val}}_{s} e^{2}}{G^{2}} \,\right],
  \qquad \vext(\bm{0}) = 0 . \;}
  \label{eq:vaspextcar}
\end{equation}

Comparing with Eq.~\eqref{eq:extpot}: the second term is \emph{exactly} the
bare-Coulomb form factor $f_s = 1$. The long-range physics was never the
problem. The first term, $v_{\mathrm{short}}$, is a cubic spline through a
\emph{tabulated} function read from the \file{POTCAR} --- and that is the entire
discrepancy.

Two further details are easy to miss and both are observable:

\begin{itemize}
  \item \textbf{Hard truncation.} For $G \ge$ \opt{PSGMA2} the \emph{whole}
    contribution is zeroed, Coulomb tail included; VASP does not extrapolate
    the table. For the gold potential \opt{PSGMAX} $=75.6\ \text{\AA}^{-1}$
    while a $128^3$ grid reaches $|\GG|_{\max}=106\ \text{\AA}^{-1}$, so the
    truncation is genuinely active and must be reproduced.
  \item \textbf{$G$ carries the $2\pi$.} \opt{LATT\_CUR\%B} is the reciprocal
    lattice \emph{without} $2\pi$, so \opt{G} is the ordinary angular
    wavevector in $\text{\AA}^{-1}$ --- the same convention as
    \pyclass{FieldGrid.get\_g\_vectors}.
\end{itemize}

\subsection{The \texorpdfstring{\file{POTCAR}}{POTCAR} table layout}

The layout of the \opt{local part} block is given by the reader in
\file{pseudo.F}:

\begin{lstlisting}[style=fortran]
READ(10,'(1X,A1)') CSEL          ! the "local part" marker line
READ(10,*) P(NTYP)%PSGMAX        ! <-- FIRST number is PSGMAX, not ZVAL
ALLOCATE(P(NTYP)%PSP(NPSPTS,5))
READ(10,*) (P(NTYP)%PSP(I,2),I=1,NPSPTS)
DO I=1,NPSPTS
    P(NTYP)%PSP(I,1)=(P(NTYP)%PSGMAX/NPSPTS)*(I-1)
ENDDO
P(NTYP)%PSCORE=P(NTYP)%PSP(1,2)
CALL SPLCOF(P(NTYP)%PSP(1,1),NPSPTS,NPSPTS,0._q)
\end{lstlisting}

with $\opt{NPSPTS}=1000$. The mesh is therefore uniform,
\begin{equation}
  q_i = \frac{\opt{PSGMAX}}{1000}\,(i-1),
  \qquad i = 1,\dots,1000,
\end{equation}
and the values are in $\mathrm{eV}\cdot\text{\AA}^3$.

\begin{pnote}
The first number after the marker is \opt{PSGMAX}, the table's maximum
wavevector --- \emph{not} the valence charge. The two can look alike in some
files, and the misreading was the secondary cause of the original discrepancy.
Confirmed against the gold \file{POTCAR}: the value is $75.5890395431569$, the
block holds $1001$ numbers ($=$ \opt{PSGMAX} $+$ exactly $1000$ samples), and
\opt{ZVAL} is $11$.
\end{pnote}

\opt{PSCORE} $= v_{\mathrm{short}}(q\to0) = 105.89\ \mathrm{eV}\cdot\text{\AA}^3$
for gold feeds the \opt{PSCENC} energy correction and does not enter the
potential, since $\vext(\GG=0)$ is zeroed.

\subsection{Result}

Implementing Eq.~\eqref{eq:vaspextcar} verbatim --- natural cubic spline on the
\opt{PSGMAX} mesh, analytic Coulomb tail, $\GG=0$ zeroed, truncation at
\opt{PSGMA2}, structure factors per species --- gives:

\begin{center}
\begin{tabular}{@{}lrrr@{}}
\toprule
Structure & grid & Gaussian model & \textbf{tabulated} \\
\midrule
struct\_000 & $128^3$ & $0.1338$ & $\mathbf{2.13\times10^{-5}}$ \\
struct\_001 & $120{\times}128{\times}128$ & $0.1303$ & $\mathbf{5.95\times10^{-5}}$ \\
struct\_002 & $128{\times}120{\times}128$ & $0.1285$ & $\mathbf{5.75\times10^{-5}}$ \\
\bottomrule
\end{tabular}
\end{center}

(relative $L^2$ against the reference \file{EXTCAR}; Pearson $r = 1.000000$ in
all three, MAE $0.34$\,meV.) The improvement is a factor of roughly $6000$.

\begin{pwarn}
The residual $\sim6\times10^{-5}$ is the spline boundary condition: VASP calls
\opt{SPLCOF(PSP(1,1),NPSPTS,NPSPTS,0.\_q)} whereas the Python implementation
uses a natural cubic spline. The difference is confined to the ends of the
table and is worth about $0.6$\,meV. It is documented rather than chased.
\end{pwarn}

\subsection{Why an analytic form factor cannot work}

For a single-species cell Eq.~\eqref{eq:vaspextcar} can be inverted to recover
the empirical form factor directly from a reference file,
\begin{equation}
  f(\GG) = -\,\frac{\Omega\,G^2\,\vext^{\mathrm{ref}}(\GG)}
                   {4\pi e^2 Z^{\mathrm{val}} S(\GG)} .
\end{equation}
Doing so for gold shows $f(G)$ decaying far more slowly than a Gaussian at low
$G$ and then \emph{oscillating about zero} at large $G$ --- the signature of a
repulsive pseudopotential core. No monotonic $e^{-G^2\sigma^2/2}$ can reproduce
that, which is why fitting $\sigma$ saturates near a relative $L^2$ of $0.13$
regardless of the value chosen. The best-fit width came out at
$\sigma^{*}=0.374\ \text{\AA}$ on all three structures independently, a
consistency that identified the limitation as structural rather than accidental.

\section{\texorpdfstring{\file{TAUCAR}}{TAUCAR}: the kinetic energy density}

\subsection{Definition}

The source comment gives
\opt{tau\_s(r) = hbar\^{}2/2m sum\_nk w\_k f\_nk |grad psi\_nk,s(r)|\^{}2},
i.e. the positive-definite gradient form of Eq.~\eqref{eq:taudef}, and
\emph{not} the Laplacian form. \opt{TAU\_PW} in \file{metagga.F} evaluates it
per Cartesian direction:

\begin{lstlisting}[style=fortran]
DO I=1,NPL
   G1=WDES%IGX(I,NK)+WDES%VKPT(1,NK)
   GC=(G1*LATT_CUR%B(IDIR,1)+G2*LATT_CUR%B(IDIR,2)+G3*LATT_CUR%B(IDIR,3))
   CFA(I)=GC*W%CPTWFP(...)*CITPI          ! i*2pi*(G+k)_IDIR * c_G
ENDDO
CALL FFTWAV(NPL,WDES%NINDPW(1,NK),CFAR(1),CFA(1),GRID)
DO I=1,GRID%RL%NP
   CKIN(I)=CKIN(I)+HSQDTM*REAL(CONJG(CFBR(I))*CFAR(I))*WEIGHT
ENDDO
\end{lstlisting}

The gradient is applied \emph{spectrally} as $i(\GG+\kk)$ on the plane-wave
coefficients, inverse-transformed to real space, and accumulated as
$|\partial_\alpha\psi|^2$ weighted by $w_k f_{nk}$, with
$\opt{HSQDTM}=\hbar^2/2m_e$.

\subsection{Storage}

\opt{CTAUC = CTAUC*LATT\_CUR\%OMEGA} followed by \opt{OUTCHG} means
\file{TAUCAR} stores $\tau\Omega$, exactly as \file{CHGCAR} stores $\rho\Omega$,
so that $\sum_\rr \file{TAUCAR}(\rr)/(N_1N_2N_3)$ is the plane-wave kinetic
energy of the cell in eV. For \opt{ISPIN=2} the two blocks are (total,
magnetisation), matching the \file{CHGCAR} layout.

\subsection{Verdict}

Poraqu\^e's handling of \file{TAUCAR} required no change: the positive-definite
definition, the $\Omega$ factor, the units, the shared grid and the
plane-wave-only content all matched already. That $\tau$ is the pseudo quantity
also explains an earlier observation --- the bound~\eqref{eq:hobound} was found
to hold at essentially every grid point (one violation in $32768$, at the single
voxel where spectral downsampling had rung $\tau$ slightly negative), because
$\tau$ and $\rho$ are pseudised \emph{consistently}. The bound is strictly
guaranteed only for all-electron quantities, so this was worth measuring rather
than assuming.

\section{What is absent from \texorpdfstring{\file{EXTCAR}}{EXTCAR}}

Worth stating explicitly, because it bounds what the first learned map can be
asked to do:

\begin{itemize}
  \item no PAW one-centre terms --- \file{EXTCAR} is a pseudo quantity;
  \item no non-local projectors --- not local functions of $\rr$, so they
    cannot appear in a volumetric file even in principle;
  \item no Hartree and no exchange--correlation, unlike \file{LOCPOT} which is
    $\vH+\vext$. \file{EXTCAR} is the bare ionic term, which is precisely the
    right input for the $\vext\mapsto\rho$ map;
  \item \textbf{no Ewald or ion--ion contribution.} That energy is a scalar
    (\opt{PSCENC}, \opt{TEWEN}), not a field; none is missing, because none
    belongs in this file.
\end{itemize}

\begin{pwarn}
Two paths remain untested. The multi-species sum in Eq.~\eqref{eq:vaspextcar}
uses a per-species \opt{PSGMAX} and is implemented as such, but the present data
set is pure gold. And \opt{ISPIN=2} writes a second block in both \file{CHGCAR}
and \file{TAUCAR} which the reader currently ignores --- correct for these
\opt{ISPIN=1} runs, a silent truncation otherwise.
\end{pwarn}
